Rivers are the main pathway through which plastic discarded on land reaches the ocean. If so, the amount of plastic at a river mouth should grow with the number of people living in the basin that feeds it. This dissertation tests that hypothesis in North and Central America, from the United States and southern Canada down to the Panama Canal, by combining three independent datasets: the GEBCO 2024 global relief grid, the GPW v4 (2020) population grid, and the plastic pellet samples of the Nurdle Patrol program in the NOAA marine microplastics database. Drainage basins are delineated with the Invasion Percolation Based Algorithm (IPBA), taking every land pixel adjacent to the ocean as an outlet; the 138 largest basins, which cover 90\% of the land area, form the sample. Cities are defined without administrative boundaries by the City Clustering Algorithm (CCA), whose two parameters, the density threshold $D^{*}$ and the linking distance $\ell$, were calibrated by Zipf's law over a $300 \times 50$ parameter sweep, solved in near-linear time by a union-find structure that activates pixels in decreasing order of density. The chosen criterion, which requires an exponent close to 1 simultaneously for the 50, 100 and 200 largest clusters, yields $D^{*} = 1010$ inhabitants/km$^2$ and $\ell = 1.7$ km, with 12,185 urban clusters and Hill exponents between 0.93 and 1.08. Clusters that cross drainage divides are split among basins. Each basin is assigned the sampling site nearest to its mouth, and the acceptance radius is swept in 1 km steps up to 50 km. At 38 km, the radius that contains most sampling sites, 52 basins have a measurement, 29 of them with a positive count. In those 29 basins, pellet concentration grows sublinearly with population, with exponent $\beta = 0.28$ (95\% CI: 0.11 to 0.54; $r = 0.37$; $p = 0.05$), and the quantity proxy, concentration times area, grows with $\beta = 0.85$ ($r = 0.66$). The association is weak and fragile: when basins with zero counts are included, the rank correlation drops to $\rho_s = 0.25$ ($p = 0.07$), and much of the exponent of the quantity proxy comes from the correlation between basin population and area. With a 1 km radius only three basins have a measurement. We conclude that current public data support, at most, a weak sublinear relation between basin population and microplastic at the mouth, and that the main limitation is the spatial coverage of the samples, concentrated in the Gulf of Mexico. The dissertation contributes a reproducible pipeline that integrates invasion percolation, Zipf-calibrated urban clustering and basin-level allometry.

\keywords{microplastics; drainage basins; invasion percolation; Zipf's law; urban scaling.}
